%% kapa: exam=1 round=35 year=2024 period=2 form=A subject=accounting subject_ko=회계학 seq=27 number=67
%% source: Q-Net 감정평가사 기출문제 (cst00309) artlSeq=5213229 file=raw/1008/5213229/2교시 문제지 원본.pdf
%% conversion: kordoc 4.24.51 -> tools/gen1.py (auto); review=proofread-sonnet; flags=none
\begin{question}{67}
무형자산의 회계처리에 관한 설명으로 옳은 것을 모두 고른 것은?
\begin{framedtext}
ㄱ. 경영자가 의도하는 방식으로 운용될 수 있으나 아직 사용하지 않고 있는 기간에 발생한 원가는 무형자산의 장부금액에 포함한다. \par
ㄴ. 자산을 사용가능한 상태로 만드는데 직접적으로 발생하는 종업원 급여와 같은 직접 관련되는 원가는 무형자산의 원가에 포함한다. \par
ㄷ. 최초에 비용으로 인식한 무형항목에 대한 지출은 그 이후에 무형자산의 원가를 신뢰성 있게 측정할 수 있다면 무형자산으로 인식할 수 있다. \par
ㄹ. 새로운 지역에서 또는 새로운 계층의 고객을 대상으로 사업을 수행하는데서 발생하는 원가 등은 무형자산 원가에 포함하지 않는다.
\end{framedtext}
\begin{choices}
  \item ㄱ, ㄴ
  \item ㄱ, ㄷ
  \item ㄱ, ㄹ
  \item ㄴ, ㄷ
  \item ㄴ, ㄹ
\end{choices}
\end{question}
\answer{5}
